% SPDX-License-Identifier: Apache-2.0
\section{An Address Map as a Trait Constant}
\label{sec:x-map}

A unit over \code{N} peripherals decodes an address against \code{N}
ranges, and stable Rust lacks array const parameters to pass them:
an \code{[usize; N]} as a const parameter needs
\code{adt\_const\_params}. Until now a router or a bridge took a
\code{BASE} and a \code{MASK} const parameter per peripheral, which is
why each count of them was a unit of its own, written out by a macro.

So the map is a type. \code{txhdl::map::AddrMap<N>} has one
associated constant, \code{RANGES: [(usize, usize); N]}, the base
and the mask of each peripheral; a design writes a marker type,
implements the trait for it with its count, and names the marker
where it names the unit, \code{Decode<3, Three>}. The unit is generic
over the marker, holds it in a \code{PhantomData} field the netlist
ignores, and reads \code{M::RANGES[i]} in a loop over \code{0..N}.
The lowering unrolls the loop when \code{lowered} runs, as
Section~\ref{sec:x-loop} unrolls one over an array of ports, and
writes each base and mask into the netlist as a number, the way it
writes a const parameter: \code{M::RANGES[i].1} is Rust worked out
at that moment, with \code{i} the unrolled loop's variable. The
array's length is checked against \code{N} where the map is written,
so a map of two ranges cannot be given to a decoder of three.

The trait has a second constant, \code{NAMES}, a few words for each
range, which a map that names nothing leaves empty. The decoder does
not read it; a table of the map does, and the run below prints each
address's range by it. The Vreteno board's maps name their ranges,
and \code{//tools/memmap} writes the address tables of the Vreteno and
flagship documents from them.

The decoder below takes an address a cycle and says, one bit per
range, which range holds it, the first that matches, and whether
none does. The run walks seven addresses in and out of the ranges;
the netlist is checked against the run under nvc and Verilator, and
its three comparators instantiate the constants defined by the map.

\begin{figure}[htbp]
\centering
\input{map_timing.tex}
\caption{The decoder over three ranges: an address a cycle, the
range that holds it as a one-hot, and \code{none} where no range
does.}
\label{fig:x-map}
\end{figure}

\lstinputlisting[caption={\code{ex\_map.rs}. Run with \code{bazel run //lib/examples:ex\_map}.},label={lst:x-map}]{lib/examples/ex_map.rs}

\lstinputlisting[style=ascii,caption={What \code{ex\_map} printed when the build ran it: the walk, then the Verilog, the map's bases and masks as numbers of the netlist.},label={lst:x-map-out}]{out_map.txt}
